\documentclass[10pt,a4paper]{amsart}
\usepackage[margin=19mm]{geometry}
\usepackage{amsmath,amssymb,mathtools}
\usepackage[scr=rsfs,cal=euler]{mathalfa}
\usepackage{xfrac}
\usepackage[unicode,hidelinks,bookmarks=false]{hyperref}
\newcommand{\ie}{i.e.\ }
\newcommand{\cf}{cf.\ }
\newcommand{\resp}{respectively\ }
\newcommand{\Subsection}{\subsection}
\providecommand{\nobreakdash}{\nobreak\hbox{-}}
\setlength{\emergencystretch}{2em}
\multlinegap0pt
\usepackage{amssymb}
\usepackage{xcolor}
\usepackage[shortlabels]{enumitem}
\usepackage{float}
\newtheorem{lemma}{Lemma}
\title{A sharp point-sphere incidence bound for $(u, s)$-Salem sets}
\author{Steven Senger, Dung The Tran}
\date{Source: arXiv:2601.07105v4 (CC BY 4.0)}
\begin{document}
\maketitle

\noindent\emph{Source and licence.} A sharp point-sphere incidence bound for $(u, s)$-Salem sets. Version arXiv:2601.07105v4 (CC BY 4.0), distributed by arXiv under the Creative Commons Attribution 4.0 International licence, http://creativecommons.org/licenses/by/4.0/, as recorded in the arXiv licence metadata for this exact version and shown on the abstract page https://arxiv.org/abs/2601.07105v4. Full licence notice: \texttt{licenses/arxiv-2601.07105-CC-BY-4.0.txt}.\par\smallskip

\noindent\textit{Editorial scope.} Counted unit: the article's lemma lem:incidence-us (source0.tex lines 1082--1109). The article's complete original proof of it is delivered with it (source0.tex lines 1111--1201, one \texttt{proof} environment of 91 lines). No mathematics is written, repaired or re-derived: the statement and the proof are the article's own lines, in the article's own notation, and no retained line is substituted or re-flowed.  No further source-local context is needed: every cross-reference of every retained line resolves inside the two delivered spans.  Nothing was omitted from any retained span, so the package carries no omission record.  No retained line contains a citation, so the wrapper carries no bibliography and no bibliography entry.  No retained line contains a figure environment, an image-inclusion command or drawing code, so no image is delivered and the package declares no asset dependency.  Editorial formatting only, in addition: an AMS \texttt{amsart} preamble that restates the article's own declarations and macros used by the retained lines, namely the environments \texttt{lemma} on separate counters, together with the standard packages those lines need.  The retained environments print in this document as Lemma 1 in order of appearance, whereas the article prints the same items with the numbering of its own full document.  The complete native source of the article as distributed in the arXiv e-print for this exact version is bundled unchanged as \texttt{sources/192474/source0.tex}.
\begin{lemma}\label{lem:incidence-us}
Let $d\ge 2$ and let $u=2k$ with $k\ge 2$.
Let $P\subset \mathbb{F}_q^{\,d}$ be a set such that
\begin{equation}\label{eq:us-salem-fourier}
\sum_{m\in\mathbb{F}_q^d\setminus\{0\}} |\widehat{P}(m)|^{u}
\ \ll_u\
q^{-(u-1)d}\,|P|^{u(1-s)},
\end{equation}
where
\[
\widehat{P}(m):=q^{-d}\sum_{x\in\mathbb{F}_q^d} P(x)\,\chi(-m\cdot x).
\]
Let $\mathcal{T}\subset (\mathbb{F}_q^d\setminus\{0\})\times \mathbb{F}_q$, and define
\[
N(P,\mathcal{T})
:=
\bigl|\{(x,(a,b))\in P\times \mathcal{T}:\ a\cdot x=b\}\bigr|.
\]
Then
\[
N(P,\mathcal{T})
\ \le\
\frac{|P|\,|\mathcal{T}|}{q}
\ +\
C_u\,|\mathcal{T}|^{1-\frac{1}{u}}\,q^{\frac{d}{u}}\,|P|^{1-s},
\]
where $C_u>0$ depends only on $u$.
\end{lemma}

\begin{proof}
For $(a,b)\in (\mathbb{F}_q^d\setminus\{0\})\times \mathbb{F}_q$, write
\[
N(a,b):=\bigl|\{x\in P:\ a\cdot x=b\}\bigr|,
\qquad
D(a,b):=N(a,b)-\frac{|P|}{q}.
\]
Then
\[
N(P,\mathcal{T})
=\frac{|P|\,|\mathcal{T}|}{q}
+
\sum_{(a,b)\in \mathcal{T}} D(a,b).
\]
By H\"older's inequality,
\begin{equation}\label{eq:holder-step-us}
\sum_{(a,b)\in \mathcal{T}} D(a,b)
\ \le\
|\mathcal{T}|^{1-\frac1u}\Bigl(\sum_{(a,b)\in \mathcal{T}} |D(a,b)|^{u}\Bigr)^{\frac{1}{u}}.
\end{equation}

We bound the $u$-th moment by enlarging the sum to all $(a,b)$ with $a\ne 0$.
By character orthogonality,
\[
N(a,b)
=
\frac{1}{q}\sum_{t\in\mathbb{F}_q}\sum_{x\in\mathbb{F}_q^d}
P(x)\,\chi\!\bigl(-t(a\cdot x-b)\bigr).
\]
Separating the term $t=0$ gives
\begin{equation}\label{eq:Dab-formula-us}
D(a,b)=q^{d-1}\sum_{t\in\mathbb{F}_q^\ast}\chi(tb)\,\widehat{P}(ta).
\end{equation}
Therefore,
\[
\sum_{(a,b)\in \mathcal{T}}|D(a,b)|^{u}
\ \le\
\sum_{a\ne 0}\sum_{b\in\mathbb{F}_q}|D(a,b)|^{u}
=
q^{u(d-1)}\sum_{a\ne 0}\sum_{b\in\mathbb{F}_q}
\Bigl|\sum_{t\ne 0}\chi(tb)\widehat{P}(ta)\Bigr|^{u}.
\]

Fix $a\ne 0$.
Expanding the $u$-th power and summing over $b\in\mathbb{F}_q$ enforces the constraint
$t_1+\cdots+t_k=t_{k+1}+\cdots+t_{2k}$.
Using the standard inequality
$c_1\cdots c_{2k}\le \frac{1}{2k}\sum\limits_{j=1}^{2k} c_j^{2k}$ for nonnegative $c_j$,
one obtains
\[
\sum_{b\in\mathbb{F}_q}\Bigl|\sum_{t\ne 0}\chi(tb)\widehat{P}(ta)\Bigr|^{u}
\ \ll_u\
q^{u-1}\sum_{t\ne 0} |\widehat{P}(ta)|^{u}.
\]
Thus,
\[
\sum_{a\ne 0}\sum_{b\in\mathbb{F}_q}|D(a,b)|^{u}
\ \ll_u\
q^{u(d-1)}\cdot q^{u-1}\sum_{a\ne 0}\sum_{t\ne 0} |\widehat{P}(ta)|^{u}.
\]
Using the change of variables $m=ta$, the map $(a,t)\mapsto m$ from $\{a\ne 0,\ t\ne 0\}$ onto
$\{m\ne 0\}$ has fiber size $q-1$. Hence
\[
\sum_{a\ne 0}\sum_{t\ne 0} |\widehat{P}(ta)|^{u}
=
(q-1)\sum_{m\ne 0}|\widehat{P}(m)|^{u}
\ll
q\sum_{m\ne 0}|\widehat{P}(m)|^{u}.
\]
Consequently,
\[
\sum_{a\ne 0}\sum_{b\in\mathbb{F}_q}|D(a,b)|^{u}
\ \ll_u\
q^{ud}\sum_{m\ne 0}|\widehat{P}(m)|^{u}.
\]
By \eqref{eq:us-salem-fourier},
\[
\sum_{a\ne 0}\sum_{b\in\mathbb{F}_q}|D(a,b)|^{u}
\ \ll_u\
q^{ud}\cdot q^{-(u-1)d}|P|^{u(1-s)}
=
q^{d}|P|^{u(1-s)}.
\]
Combining this with \eqref{eq:holder-step-us} yields
\[
\sum_{(a,b)\in \mathcal{T}} D(a,b)
\ \ll_u\
|\mathcal{T}|^{1-\frac1u}\,q^{\frac{d}{u}}\,|P|^{1-s}.
\]
This proves the lemma.
\end{proof}

\end{document}
